% ===== PRÉAMBULE CANONIQUE — à coller VERBATIM en tête de CHAQUE .tex =====
\documentclass[11pt,a4paper]{article}
\usepackage[utf8]{inputenc}\usepackage[T1]{fontenc}\usepackage{lmodern}
\usepackage[french]{babel}
\usepackage{amsmath,amssymb,mathtools}\usepackage{icomma}
\usepackage[version=4]{mhchem}          % \ce{...} : équations & formules
\usepackage{siunitx}\sisetup{locale=FR} % \qty{}{}, \si{}, \ang{}
\usepackage{chemfig}                    % \chemfig{...} : structures développées
\usepackage{xcolor}\usepackage{colortbl}
\usepackage{tikz}\usepackage{pgfplots}\pgfplotsset{compat=1.18}
\usepackage[most]{tcolorbox}\usepackage{enumitem}\usepackage{array,booktabs}
\usepackage[a4paper,margin=2.2cm]{geometry}
\usetikzlibrary{arrows.meta,calc,angles,quotes,positioning,patterns,backgrounds,shapes.geometric}
\definecolor{primary}{RGB}{222,49,99}\definecolor{primarydark}{RGB}{150,22,55}
\definecolor{defcol}{RGB}{0,114,178}\definecolor{methcol}{RGB}{52,92,148}
\definecolor{excol}{RGB}{0,135,90}\definecolor{attncol}{RGB}{200,40,55}
\tcbset{boxbase/.style={enhanced,breakable,boxrule=0.9pt,arc=2.5pt,
  left=3mm,right=3mm,top=2.3mm,bottom=2.3mm,fonttitle=\bfseries,coltitle=white}}
% --- boîtes TUTORIEL (noms EXACTS, ne pas renommer) ---
\newtcolorbox{cle}[1][]{boxbase,colback=defcol!6,colframe=defcol,
  title={\textsc{À retenir}\ifx\relax#1\relax\relax\else\textmd{ : \itshape #1}\fi}}
\newtcolorbox{methode}[1][]{boxbase,colback=methcol!7,colframe=methcol,sharp corners,
  title={\textsc{Méthode}\ifx\relax#1\relax\relax\else\textmd{ --- \itshape #1}\fi}}
\newtcolorbox{exemplebox}[1][]{boxbase,colback=excol!5,colframe=excol,
  title={\textsc{Exemple}\ifx\relax#1\relax\relax\else\textmd{ : \itshape #1}\fi}}
\newtcolorbox{piege}[1][]{boxbase,colback=attncol!6,colframe=attncol,
  title={\textsc{Piège}\ifx\relax#1\relax\relax\else\textmd{ : \itshape #1}\fi}}
\newtcolorbox{remarque}[1][]{boxbase,colback=black!4,colframe=black!45,
  title={\textsc{Remarque}\ifx\relax#1\relax\relax\else\textmd{ : \itshape #1}\fi}}
% --- boîtes EXERCICE / CORRIGÉ (noms EXACTS, auto-comptées) ---
\newtcolorbox[auto counter]{exo}[1][]{boxbase,colback=primary!4,colframe=primary,
  title={\textsc{Exercice}~\thetcbcounter\ifx\relax#1\relax\relax\else\textmd{ --- \itshape #1}\fi}}
\newtcolorbox[auto counter]{corrige}[1][]{boxbase,colback=excol!4,colframe=excol,
  title={\textsc{Corrigé}~\thetcbcounter\ifx\relax#1\relax\relax\else\textmd{ --- \itshape #1}\fi}}
\newcommand{\R}{\mathbb{R}}\newcommand{\N}{\mathbb{N}}\newcommand{\Z}{\mathbb{Z}}\newcommand{\ds}{\displaystyle}
% =========================================================================
\begin{document}
\sisetup{per-mode=symbol}
\begin{exo}[calculer le quotient réactionnel $Q$]
Écris l'expression de $Q$ puis calcule sa valeur avec les concentrations effectives données.
\begin{enumerate}[label=\alph*)]
  \item \ce{AgCl} : $[\ce{Ag+}]=\qty{1.0e-4}{\mole\per\litre}$, $[\ce{Cl-}]=\qty{1.0e-5}{\mole\per\litre}$.
  \item \ce{PbBr2} : $[\ce{Pb^{2+}}]=\qty{1.0e-3}{\mole\per\litre}$, $[\ce{Br-}]=\qty{1.0e-2}{\mole\per\litre}$.
  \item \ce{Mg(OH)2} : $[\ce{Mg^{2+}}]=\qty{1.0e-2}{\mole\per\litre}$, $[\ce{OH-}]=\qty{1.0e-4}{\mole\per\litre}$.
\end{enumerate}
\end{exo}

\begin{exo}[comparer $Q$ et $K_{\mathrm{ps}}$]
On donne $K_{\mathrm{ps}}(\ce{AgCl})=\num{1.8e-10}$. Pour chaque valeur de $Q$, indique l'état de la
solution (non saturée / saturée / sursaturée) et s'il y a précipitation.
\begin{enumerate}[label=\alph*)]
  \item $Q=\num{1.0e-9}$
  \item $Q=\num{1.0e-12}$
  \item $Q=\num{1.8e-10}$
\end{enumerate}
\end{exo}

\begin{exo}[une solution de \ce{SnS} est-elle saturée ?]
Le sulfure d'étain \ce{SnS} (type $1{:}1$) a $K_{\mathrm{ps}}=\num{1.0e-26}$. Une solution contient
$[\ce{Sn^{2+}}]=\qty{1.0e-19}{\mole\per\litre}$.
\begin{enumerate}[label=\alph*)]
  \item Calcule la solubilité à saturation $s_{\text{sat}}=\sqrt{K_{\mathrm{ps}}}$.
  \item Compare $[\ce{Sn^{2+}}]$ à $s_{\text{sat}}$ : la solution est-elle saturée ?
\end{enumerate}
\end{exo}

\begin{exo}[effet d'ion commun : $s'=K_{\mathrm{ps}}/C^{n}$]
Calcule la nouvelle solubilité $s'$ du sel en présence de l'ion commun à la concentration $C$ imposée.
\begin{enumerate}[label=\alph*)]
  \item \ce{AgCl} ($K_{\mathrm{ps}}=\num{1.8e-10}$) dans $[\ce{Cl-}]=\qty{0.10}{\mole\per\litre}$.
  \item \ce{AgCl} ($K_{\mathrm{ps}}=\num{1.8e-10}$) dans $[\ce{Cl-}]=\qty{0.010}{\mole\per\litre}$.
  \item \ce{Mg(OH)2} ($K_{\mathrm{ps}}=\num{9.0e-12}$) dans $[\ce{OH-}]=\qty{0.10}{\mole\per\litre}$.
\end{enumerate}
\end{exo}

\begin{exo}[vrai ou faux : précipitation et ion commun]
Indique si chaque affirmation est vraie ou fausse.
\begin{enumerate}[label=\alph*)]
  \item L'ajout d'un ion commun \emph{augmente} la solubilité du sel.
  \item L'ajout d'un ion commun modifie la valeur de $K_{\mathrm{ps}}$.
  \item Si $Q>K_{\mathrm{ps}}$, un précipité se forme.
  \item Dans une solution saturée à l'équilibre, $Q=K_{\mathrm{ps}}$.
\end{enumerate}
\end{exo}
\end{document}
